\subsection{Relation to Q-ball excitations and continuum trapping}
\label{sec:related-work}

Q-balls are localized scalar configurations carrying a conserved global
charge~\cite{Coleman1985}. Battye and Sutcliffe study collisions, charge
transfer and fission in one, two and three spatial dimensions in the same
sextic model, after unit conversion~\cite{BattyeSutcliffe2000}.
Kovtun, Nugaev, and Shkerin formulate the coupled sideband problem
$\omega\pm\rho$ and analyze three-dimensional vibrational
modes~\cite{Kovtun2018}. Azatov, Ho, and Khalil treat scattering in the
sextic family, as well as the complex-plus-real-field
Friedberg--Lee--Sirlin model~\cite{Azatov2024,FriedbergLeeSirlin1976}.
The coupled radial operator is established prior work; the contribution
here concerns particular localized solutions of it, with regular origin
conditions and controlled infinite tails.

Chen, Andersson, and Li analyze three-dimensional Q-ball perturbations
for the rescaled potential $V(s)=s^2-s^3+\kappa s^4$, where $s$ denotes
the field modulus~\cite{ChenAnderssonLi2025}. Their coupled-sideband
classification provides context, but concerns a different potential.
Our claim is the certification of three particular localized linear
modes of the sextic model, including cancellation of both oscillatory
open-channel components. Among the works reviewed here, we did not
identify a predecessor for these specific modes. This limited
statement establishes neither global priority nor a contradiction with
the spectrum studied by Chen and collaborators.

Zhang, Zhou, and Zhu obtain analytic scattering solutions from stepwise
Q-ball profiles~\cite{ZhangZhouZhu2025}. Their interior dispersion and
thin-wall phase structure are directly relevant to Eq.~\eqref{eq:interior},
although their propagating scattering regime has two exterior open
channels. Ciurla et al. distinguish bound, half-propagating and quasinormal
modes in $1+1$ dimensions~\cite{Ciurla2024}. Their half-propagating interval
$1-\omega<\rho<1+\omega$ has the same channel thresholds as ours, but their
long-lived example retains a small radiative tail. Evslin et al. analyze
paired Floquet modes of small-amplitude, one-dimensional
Q-balls~\cite{Evslin2026}; coupling of a discrete component to its
propagating partner produces Feshbach-type quasinormal modes. These
one-dimensional calculations do not transfer the spherical origin
condition or radial measure to three dimensions.

Numerically assisted spectral exclusion arguments for NLS solitary waves
include those of Marzuola and Simpson and Asad and
Simpson~\cite{MarzuolaSimpson2011,AsadSimpson2011}.
Cuccagna and Maeda impose absence of embedded eigenvalues as hypothesis H4
and describe the broader ground-state expectation as unproved in their
2020 survey~\cite[Remark 3.2]{CuccagnaMaeda2020}.
Li and Yang subsequently prove absence in the essential spectrum for the
full linearized operator about the ground state of the three-dimensional
focusing cubic NLS~\cite[Theorem 1.1]{LiYang2026}. This is a different
equation and operator from our relativistic Q-ball frequency pencil.
A nodeless Q-ball profile alone does not establish the relevant
variational ground-state class. The present model-specific existence
claim therefore does not contradict that NLS theorem or refute the
survey's expectation in its stated setting.

A bound state in the continuum (BIC) is a spatially localized linear state
at a frequency also admitting extended radiation modes~\cite{Hsu2016}.
Interference-induced trapping is an established mechanism:
Friedrich and Wintgen relate interacting resonances to vanishing decay
width~\cite{FriedrichWintgen1985}, and Yu and Lu give an existence theorem
for a three-channel Schr\"odinger problem with one open
channel~\cite{YuLu2025}. Its small-coupling assumptions and finite-range
potential differ from the present two-channel pencil with an infinite
profile tail. Related resonant reflection occurs for discrete breathers,
optical solitons and condensate domain
walls~\cite{Flach2003,Flach2005,Watabe2012}.

Winding numbers for BICs are also established: Zhen et al. use the winding
of a real far-field polarization vector in the two-dimensional Bloch
wave-vector plane of a photonic slab~\cite{Zhen2014}. Their symmetry
assumptions support a charge-conservation and robustness statement.
Our winding of a matching map over Q-ball parameters is a numerical
zero-finding diagnostic; it does not inherit that photonic robustness
theorem.

The planar transmission zero and interior phase matching admit a
\emph{Fabry--P\'erot interpretation}. Hsu et al. describe BICs formed by
perfect resonant reflection together with a round-trip phase condition,
including a resonator paired with a reflecting boundary~\cite{Hsu2016}.
In the thin-wall picture used here, the reflecting wall and the regular
centre supply the analogous radial cavity. This is an interpretation of
the wall and phase calculations, rather than an additional existence
result for the curved Q-ball. In particular, the planar zero alone does
not certify a spherical mode. Sequences of radiation cancellations and
WKB matching are also known for nonlinear embedded
solitons~\cite{Malomed2005}; their model and scaling parameter differ from
those of the present sequence of linear modes on changing backgrounds.

Nonlinear radiation obstructions have a separate history. Segur and
Kruskal find beyond-all-orders radiation excluding small-amplitude true
breathers in the one-dimensional $\phi^4$
model~\cite{SegurKruskal1987}. Sigal studies instability of periodic and
quasiperiodic wave solutions under generic nonlinear
perturbations~\cite{Sigal1993}; Soffer and Weinstein prove radiation
damping under a Fermi-golden-rule hypothesis~\cite{SofferWeinstein1999}.
Manton and Merabet derive second-harmonic radiation with amplitude of
order $A^2$ and power of order $A^4$ from a kink's internal mode in
$1+1$-dimensional $\phi^4$ theory~\cite[Appendix A]{MantonMerabet1997}.
The harmonic mechanism and power counting are known; their coefficients
and perturbative decay law are not transferred to our three-dimensional
Q-ball.
The present contribution in this direction is the explicit error budget
separating a particular second-harmonic coefficient from zero. Its
continuation obstruction has the stated hypotheses and is not a
dynamical decay theorem. Oscillon decay minima~\cite{ZhangOscillons2020}
and the second-harmonic radiation zeros of a critical-bubble internal
mode~\cite{InagakiMurakami2026} further illustrate why suppression at one
harmonic does not establish an exactly localized nonlinear oscillation.

L\"uhrmann et al. combine analytic estimates and interval arithmetic to
certify vortex profiles, spectral data and nonzero Fermi-golden-rule
coefficients in the self-dual abelian Yang--Mills--Higgs model in two
spatial dimensions~\cite[Theorem 1.1, Remark 1.2 and Section 2]{LuhrmannEtAl2026}.
Their internal mode is a gap eigenvalue. This is a methodological
precedent for rigorous spectral and radiation-coupling computations.
Here the distinct target is a particular outgoing second-harmonic
coefficient associated with a localized linear Q-ball mode in three
dimensions. Its conditional continuation obstruction remains separate
from a nonlinear decay, lifetime or background-stability theorem.

Background stability is likewise distinct from mode localization.
Classical stability results for charged scalar solitons and the
Vakhitov--Kolokolov slope criterion provide the appropriate
context~\cite{FriedbergLeeSirlin1976,VakhitovKolokolov1973}.
Applying such a criterion requires its model-dependent hypotheses and
charge convention; the certificates here do not determine $dQ/d\omega$
at the three backgrounds. A numerical evaluation and its limited scope
are given in Section~\ref{sec:discussion}. The literature comparisons delimit the claimed
contribution without establishing exhaustive priority or an infinite
ladder.
